% The real-or-modelled map: what a joint in a robot has, and what stands in for
% it on this bench. This figure is the chapter's data, and it is printed on the
% console at every boot as one line of text.
\begin{tikzpicture}[font=\sffamily\small, node distance=3mm]

\node[font=\sffamily\small\bfseries] (h1) at (0,0) {a joint in a robot};
\node[font=\sffamily\small\bfseries] (h2) at (7.4,0) {this bench, at chapter 20};

% ---------------- left column: the real thing
\node[hw, below=6mm of h1.south, text width=32mm] (r1) {position sensing};
\node[hw, below=of r1, text width=32mm] (r2) {rate and acceleration};
\node[hw, below=of r2, text width=32mm] (r3) {periodic control loop};
\node[hw, below=of r3, text width=32mm] (r4) {motor and drive stage};
\node[hw, below=of r4, text width=32mm] (r5) {torque measurement};
\node[hw, below=of r5, text width=32mm] (r6) {field bus to the master};
\node[hw, below=of r6, text width=32mm] (r7) {real-time Ethernet slave};
\node[hw, below=of r7, text width=32mm] (r8) {certified safe stop};

% ---------------- right column: what this bench has
\node[modelled,  right=42mm of r1, text width=36mm] (b1) {half real: the decode path is hardware, the source is generated};
\node[fw,        right=42mm of r2, text width=36mm] (b2) {real: inertial unit on a shield};
\node[fw,        right=42mm of r3, text width=36mm] (b3) {real: timer driven, jitter measured};
\node[modelled,  right=42mm of r4, text width=36mm] (b4) {half real: the outputs and dead time are real, the plant is a model};
\node[modelled,  right=42mm of r5, text width=36mm] (b5) {modelled: derived from the plant's own output};
\node[fw,        right=42mm of r6, text width=36mm] (b6) {real: CAN-FD on a real transceiver};
\node[absentblk, right=42mm of r7, text width=36mm] (b7) {absent: no Ethernet controller, no slave chip};
\node[absentblk, right=42mm of r8, text width=36mm] (b8) {absent: a safe state is implemented, nothing is certified};

% ---------------- the correspondence
\foreach \i in {2,3,6} { \draw[flow] (r\i) -- node[lbl]{is} (b\i); }
\foreach \i in {1,4,5}   { \draw[hiz, ->] (r\i) -- node[lbl]{stands in for} (b\i); }
\foreach \i in {7,8}     { \draw[dotted, draw=black!55, ->] (r\i) -- node[lbl]{explained} (b\i); }

% ---------------- the chapter that settles each row
\foreach \i/\c in {1/5, 2/6, 3/2, 4/7, 5/8, 6/9, 7/17, 8/18}
  { \node[chlabel, anchor=west] at ($(b\i.east)+(1.5mm,0)$) {ch \c}; }

\node[umlnote, text width=74mm, below=7mm of b8.south east, anchor=north east]
  {Three rows are real, three stand in, two are absent and explained. The eight
   rows are the file \texttt{doc/real-or-modelled.md}, they are the line the
   board prints after its name at every boot, and they are the honest answer to
   whether this node has been built. A demonstration that blurs any row of this
   figure has failed, however well it runs.};
\end{tikzpicture}
